%قبل از شروع مسئله ترجیحاً چیزی ننویسید. اما اگر خواستید بنویسید، در قالب لیستی (فقرات و شمارش) بنویسید.
%برای اضافه کردن نمره، بعدی از شروع‌مسئله، [x] را اضافه کنید و به جای x نمره مورد نظر را قرار دهید. برای زیرمسئله‌ها هم به همین صورت است.
%\شروع‌مسئله[10]
\شروع‌مسئله [15]
%\متن‌سیاه{سوال عملی اول:}\\
%\متن‌سیاه{دراپ‌تست}\\
%\متن‌سیاه{ورودی:}\\
%در خط اول ورودی، ابتدا k و سپس m داده می‌شود.\\
%\متن‌سیاه{خروجی:}\\
%در تنها سطر خروجی، تعداد آیفون‌های مورد نیاز می‌آید. اگر در هیچ صورتی چنین کاری ممکن نبود ۱- چاپ می‌شود.\\
%\متن‌سیاه{محدودیت‌ها:}\\
%$1\leq k\leq 10^4$

%$m\leq10^9$
%\\
%\textbf{\متن‌سیاه{ورودی و خروجی نمونه}\\
%% تست اول
%\begin{latin}
% \begin{tabularx}{\textwidth}{|X | X|}
% \hline
% \textbf{stdin} &
% \textbf{stdout}\\
% \hline
% 100 100&
% 1\\
% \hline
% 
% \end{tabularx}
%\end{latin}
%
%% تست دوم
%
%\begin{latin}
% \begin{tabularx}{\textwidth}{|X | X|}
% \hline
% \textbf{stdin} &
% \textbf{stdout}\\
% \hline
% 340 20&
% 3\\
% \hline
% 
% \end{tabularx}
%\end{latin}}
% صورت اصلی مسئله را در این‌جا بنویسید.

الی و جول در \textit{آخرین ما}\پانویس{The Last of Us} از اکتشاف در دنیای آخرالزمانی‌شان خسته شده‌اند و تصمیم می‌گیرند بازی جدیدی برای وقت‌گذرانی انجام دهند.

آن‌ها برای این کار دنبالۀ اعداد طبیعی دلخواه $a_1,a_2,\dots ,a_n$ را در نظر می‌گیرند. به ازای هر جفت $a_i$ و $a_j$ $(1\le i<j\le n)$، اگر $a_i>a_j$ باشد الی $a_i+a_j$ امتیاز، اگر $a_i<a_j$، جول $a_i+a_j$ امتیاز و اگر $a_i=a_j$ هر کدام $(a_i + a_j) / 2$ امتیاز می‌گیرند. کسی که در نهایت امتیاز بیشتری داشته باشد، برنده است. 

پس از مدتی، آن‌ها از این بازی هم خسته می‌‌شوند و می‌خواهند بدانند چه کسی برنده می‌شود.

\متن‌سیاه{ورودی:}\\
در خط اول ورودی عدد $n$ داده می‌شود که تعداد اعداد دنباله است. در خط بعدی $n$ عدد دنباله با فاصله از هم می‌آیند.

\متن‌سیاه{خروجی:}\\
در خط اول خروجی بسته به نتیجۀ بازی \متن‌لاتین{"Ellie"}، \متن‌لاتین{"Joel"} یا \متن‌لاتین{"draw"} (بدون نقل قول)، و در خط دوم امتیاز برنده (در صورت تساوی امتیاز هر کدام) را نمایش دهید.

\متن‌سیاه{محدودیت‌ها:}\\
$1\leq n\leq 3 \times 10^5$\\
$1\leq a_i\leq 10^{8}$

\متن‌سیاه{ورودی و خروجی نمونه:}\\


%% تست اول
\begin{latin}
	\begin{tabularx}{\textwidth}{|X | X|}
		\hline
		\textbf{stdin} &
		\textbf{stdout}\\
		\hline
		3
		
		7 1 1&
		Ellie
		
		17\\
		\hline
		
	\end{tabularx}
\end{latin}


\begin{latin}
	\begin{tabularx}{\textwidth}{|X | X|}
		\hline
		\textbf{stdin} &
		\textbf{stdout}\\
		\hline
		4
		
		3 1 4 2&
		draw
		
		15\\
		\hline
		
	\end{tabularx}
\end{latin}

\begin{latin}
	\begin{tabularx}{\textwidth}{|X | X|}
		\hline
		\textbf{stdin} &
		\textbf{stdout}\\
		\hline
		5
		
		1 2 3 4 5&
		Joel
		
		60\\
		\hline
		
	\end{tabularx}
\end{latin}

%% تست دوم

%اگر مسئله چند قسمت داشت به این صورت عمل کنید:
%\شروع{زیرمسائل}
%\زیرمسئله[0] % مسئلۀ ۱آ
%\زیرمسئله % مسئلۀ ۱ب
%\زیرمسئله % مسئلۀ ۱پ
%\زیرمسئله % مسئلۀ ۱ت
%\پایان{زیرمسائل}

\solution{
	اگر $a_i$ و $a_j$ برابر باشند نصف امتیاز، و اگر در دنباله نابجایی داشته باشند (یعنی $a_i>a_j$ در حالی که $i<j$)، کل امتیاز این جفت به الی می‌رسد.
	ساده‌ترین الگوریتم، برای هر جفت $a_i$ و $a_j$ امتیاز را محاسبه می‌کند. اما پیچیدگی زمانی آن از مرتبۀ $O(n^2)$ است. با استفاده از ایده مرتب‌سازی ادغامی می‌توان الگوریتمی سریع‌تر با مرتبۀ زمانی $O(n\log n)$ ارائه داد.
	
	الگوریتم مرتب‌سازی ادغامی را برای دنباله اجرا کنید. مرحله ادغام را در این الگوریتم در نظر بگیرید. هر دو عدد $a_i$ و $a_j$ که نابجایی داشته باشند ($a_i>a_j$ و $i<j$)، دقیقا یک بار در دو دنبالۀ مختلف قرار می‌گیرند که الگوریتم می‌خواهد ادغامشان کند. این دو دنباله را $b_1\le b_2\le \dots \le b_m$ و $c_1\le c_2\le \dots \le c_k$ بگیرید که $a_i$ در دنبالۀ اول و $a_j$ در دنبالۀ دوم قرار دارد؛ یعنی $t$‌ای وجود دارد که $c_t = a_j$. الگوریتم ادغام اشاره‌گر‌هایی را در ابتدای هر دو دنباله قرار می‌دهد و هر کدام را که عدد کوچکتری نشان می‌داد، جلو می‌برد. این اشاره‌گر پس از زمانی به $c_t$ می‌رسد و $c_t$ را که در دنبالۀ اولیه قبل از هر کدام از اعداد $b_s, b_{s+1}, \dots ,b_m$ قرار داشت ($s$ و $t$ به ترتیب اشاره‌گر‌های اول و دوم هستند)، جلوی آن‌ها قرار می‌دهد و نابجایی‌اش را با آن‌ها را از بین می‌برد. دقت کنید که کل الگوریتم مرتب‌سازی ادغامی صرفا متشکل از چندین مرحله ادغام دو زیر دنباله از دنباله اصلی است و هر جفت نابجایی دقیقا در یکی از این ادغام‌ها ترتیب‌شان عوض می‌شود. در هر مرحله ادغام، امتیازی که الی بدست می‌آورد را محاسبه می‌کنیم. یعنی وقتی اشاره‌گر دوم به $c_t$ رسید، $(b_s + c_t) + (b_{s+1} + c_t) + \dots + (b_m + c_t) = b_s + b_{s+1} + \dots + b_m + (m - s + 1) \times c_t$ را به الی اضافه کنیم. برای اینکه این محاسبه، مرتبۀ پیچیدگی زمانی ادغام را از $O(m + k)$ بیشتر نکند، می‌توانیم در ابتدای هر مرحله ادغام مجموع‌های جزئی $b_m, b_m + b_{m-1}, \dots , b_m + b_{m-1} + \dots + b_1$ را در جایی نگه داریم. 
	
	در آخر توجه کنید که مجموع امتیاز‌های الی و جول خیلی ساده به‌دست می‌آید و کافیست فقط امتیاز یکی را به‌دست بیاوریم.
	
}




%%%%%%%%%%%%%%%%%%%%%%%%%%%%%%%%%%%%%%%%%%%%%%%%%%%%%
% برخی دستورات مفید برای تمرین‌ها. %
%%%%%%%%%%%%%%%%%%%%%%%%%%%%%%%%%%%%%%%%%%%%%%%%%%%%%

% فرمول‌های نمونه:
%$f(x) = \Theta(x)$
%$T(x, y) \leq \log(x) + 2^y + \binom{2n}{n}$
% {\tt code\_function}


% لیست موارد بدون شماره‌گذاری:
%\شروع{فقرات}
% \فقره مورد اول 
% \شروع{فقرات}
% \فقره مورد اول داخلی 
% \شروع{فقرات}
% \فقره اولین مورد خیلی داخلی 
% \فقره دومین مورد خیلی داخلی 
% \پایان{فقرات}
% \فقره مورد دوم داخلی 
% \پایان{فقرات}
% \فقره مورد دوم 
%\پایان{فقرات}

% لیست موارد شماره‌گذاری شده:
%\شروع{شمارش}
% \فقره یک مورد 
% \فقره یک مورد دیگر 
% \فقره باز هم یک مورد دیگر
%\پایان{شمارش}

% برای نوشتن یک فرمول چند خطی که یک شماره بخوره. برای حالت‌های دیگر به سایت زیر مراجعه کنید:
% https://www.overleaf.com/learn/latex/Aligning_equations_with_amsmath
%\begin{align} 
% \begin{split}
% (x+y)^3 &= (x+y)^2(x+y) \\
% &= (x^2+2xy+y^2)(x+y) \\
% &= (x^3+2x^2y+xy^2) + (x^2y+2xy^2+y^3) \\
% &= x^3+3x^2y+3xy^2+y^3
% \end{split} 
%\end{align}

% نوشتن ماتریس و آرایه دوبعدی
%\begin{align}
% A = 
% \begin{bmatrix}
% A_{11} & A_{21} \\
% A_{21} & A_{22}
% \end{bmatrix}
%\end{align}

% اضافه کردن تصویر با شرح. برای اضافه کردن تصویر بدون شرح، فقط دستور درج‌تصویر را استفاده کنید.
%\شروع{شکل}[h!]
%\تنظیم‌ازوسط
% \درج‌تصویر[پهنا=0.5\پهنای‌متن]{img.jpg}
%\شرح{شرح تصویر}
%\برچسب{تصویر:۱}
%\پایان{شکل}
%\پایان{زیرمسائل}

%کد زیر را ببینید.
%%\lstinputlisting{source_filename.py}
%\begin{latin}
%\begin{lstlisting}
%# Import the modules
%import sys
%import random
%
%ans = True
%
%while ans:
%question = raw_input("Ask the magic 8 ball a question: (press enter to quit) ")
%
%answers = random.randint(1,8)
%
%if question == "":
%sys.exit()
%
%elif answers == 1:
%print "It is certain"
%
%elif answers == 2:
%print "Outlook good"
%
%elif answers == 3:
%print "You may rely on it"
%
%elif answers == 4:
%print "Ask again later"
%
%elif answers == 5:
%print "Concentrate and ask again"
%
%elif answers == 6:
%print "Reply hazy, try again"
%
%elif answers == 7:
%print "My reply is no"
%
%elif answers == 8:
%print "My sources say no"
%\end{lstlisting}
%\end{latin}